\uuid{GjYK}
\exo7id{4183}
\titre{exo7 4183}
\auteur{quercia}
\organisation{exo7}
\datecreate{2010-03-11}
\isIndication{false}
\isCorrection{false}
\chapitre{Fonction de plusieurs variables}
\sousChapitre{Dérivée partielle}
\cpge{mpsi-mp2i.fonctions-deux-variables.b, pcsi.fonctions-deux-variables.b, ptsi.fonctions-deux-variables.b, mpsi-mp2i.integration.c, pcsi.integration.b, ptsi.integration.b, bcpst1.fonctions-deux-variables.c, bcpst1.integration-segment.d, pc.integration.a, psi.integration.a}
\module{Analyse}
\niveau{L2}
\difficulte{}

\contenu{
\question{
On pose~:
$f(x,y) = \begin{cases}
 x            &\text{ si } y\ge0 \text{ et } 0\le x\le \sqrt y~;\cr
 2\sqrt y - x &\text{ si } y\ge0 \text{ et } \sqrt y < x\le 2\sqrt y~;\cr
 0            &\text{ si } y\ge0 \text{ et } 2\sqrt y < x \text{ ou } x\le 0~;\cr
 -f(x,-y)     &\text{ si } y<0.\cr\end{cases}
$
et~: $F(y) =  \int_{x=0}^1 f(x,y)\,d x$.

Faire un dessin, vérifier que $f$ est continue sur~$\R^2$,
calculer $F(y)$ pour $-\frac14 \le y \le \frac 14$, $F'(0)$ et
$ \int_{x=0}^1 \frac{\partial f}{\partial y}(x,0)\,d x$.
}
}
